\section{Background}
\label{sec:background}

\subsection{DDoS attacks in the dataset}
CIC-DDoS2019 groups attacks into \emph{reflection-based} attacks, where third-party servers
answer spoofed requests and so hide and amplify the attacker, and \emph{exploitation-based}
attacks, which abuse protocol behaviour directly~\cite{sharafaldin2019ddos}. The five attack
classes in this project's data file are summarised in Table~\ref{tab:attacks}.

\begin{table}[H]
\centering\small
\caption{Attack classes present in \code{output.csv}.}
\label{tab:attacks}
\begin{tabularx}{\linewidth}{l l X}
\toprule
Label & Type & Mechanism \\
\midrule
\code{MSSQL} & reflection (TCP/UDP) & Spoofed requests to the SQL Server Resolution Protocol
  (UDP 1434). Servers answer with much larger instance listings. \\
\code{LDAP} & reflection (UDP) & Connectionless LDAP (CLDAP, UDP 389) queries whose responses are
  tens of times larger than the request. \\
\code{Portmap} & reflection (UDP) & Abuse of the RPC port mapper (rpcbind, UDP 111), which returns
  long service lists. \\
\code{UDP} & exploitation & High-rate flood of UDP datagrams to random ports. The victim spends
  resources answering with ICMP ``port unreachable''. \\
\code{Syn} & exploitation (TCP) & SYN flood: many half-open TCP handshakes exhaust the victim's
  connection table. \\
\code{BENIGN} & --- & Normal background traffic generated by profiled user agents. \\
\bottomrule
\end{tabularx}
\end{table}

These classes leave different traces at the flow level. SYN floods produce very short flows with
the SYN flag set and almost no backward traffic. Reflection attacks produce unidirectional flows
with characteristic payload sizes. Benign traffic is far more varied. This is why flow statistics
separate benign from attack traffic so well, and also why the \emph{reflection} classes (MSSQL,
LDAP, Portmap) are the hard part of the multiclass problem: they look alike.

\subsection{Flow-based detection and CICFlowMeter}
A \emph{flow} is the set of packets sharing a 5-tuple (source/destination IP and port, protocol)
within a timeout. CICFlowMeter~\cite{sharafaldin2018} turns packet captures into bidirectional
flows and computes about 80 statistics per flow: durations, forward/backward packet counts and
byte totals, packet-length moments, inter-arrival-time (IAT) moments, TCP flag counts, header
lengths, bulk and sub-flow statistics, and active/idle periods. Two properties of these features
matter a great deal later:
\begin{itemize}
\item \textbf{Heavy tails.} Rates such as \code{Flow Bytes/s} span nine orders of magnitude
      (Section~\ref{sec:dataset}).
\item \textbf{Redundancy.} Many columns are deterministic functions of others (e.g.\
      \code{Fwd Header Length.1} duplicates \code{Fwd Header Length}; \code{Packet Length Variance}
      is the square of \code{Packet Length Std}), and several are always zero.
\end{itemize}

\subsection{Support vector machines and kernels}
Given training data $\{(\bm{x}_i,y_i)\}_{i=1}^N$ with $y_i\in\{-1,+1\}$, a soft-margin SVM~\cite{cortes1995}
solves the dual problem
\begin{equation}
\max_{\bm{\beta}}\;\sum_{i}\beta_i-\frac12\sum_{i,j}\beta_i\beta_j y_iy_j K(\bm{x}_i,\bm{x}_j)
\quad\text{s.t.}\quad 0\le\beta_i\le C,\;\; \sum_i\beta_iy_i=0,
\label{eq:svmdual}
\end{equation}
and predicts $f(\bm{x})=\operatorname{sign}\big(\sum_i\beta_iy_iK(\bm{x}_i,\bm{x})+b\big)$.
The data enter \emph{only} through the kernel $K(\bm{x},\bm{x}')=\langle\varphi(\bm{x}),\varphi(\bm{x}')\rangle$,
the inner product in some feature space. This is the \emph{kernel trick}: any positive
semi-definite (PSD) function can be plugged in, including one evaluated on a quantum computer.
scikit-learn's \code{SVC(kernel='precomputed')} accepts the $N\times N$ Gram matrix
$K_{ij}=K(\bm{x}_i,\bm{x}_j)$ for training and the $M\times N$ matrix
$K(\bm{x}^{\text{test}}_m,\bm{x}_i)$ for prediction~\cite{sklearn}.

\paragraph{The RBF kernel.} The classical baseline is $K_{\text{RBF}}(\bm{x},\bm{x}')=\exp(-\gamma\|\bm{x}-\bm{x}'\|^2)$.
The \emph{bandwidth} $\gamma$ sets the scale at which two points count as similar. If $\gamma$ is
too large, $K\to I$ and the SVM memorises the training set. If it is too small, $K\to\bm{1}\bm{1}^\top$
and every point looks the same. This idea carries over directly to quantum kernels
(Section~\ref{sec:conc}).

\paragraph{Regularisation $C$.} $C$ trades margin width against training errors. $C$ and $\gamma$
interact strongly, which is why they are normally tuned jointly by cross-validation.

\paragraph{Multiclass and probabilities.} For $k$ classes, \code{SVC} trains $k(k-1)/2$
one-vs-one classifiers and takes a vote. Probabilities (\code{predict\_proba}) come from Platt
scaling~\cite{platt1999}, a logistic model fitted on internal cross-validation folds. Platt
scaling is randomised, so it must be seeded to be reproducible.

\subsection{Quantum computing essentials}
A qubit is a unit vector $\ket{\psi}=a\ket0+b\ket1\in\mathbb{C}^2$. An $n$-qubit register lives in
$(\mathbb{C}^2)^{\otimes n}\cong\mathbb{C}^{2^n}$, with basis states $\ket{z}$ indexed by bit-strings
$z\in\{0,1\}^n$. Gates are unitary matrices. The ones used here are
\begin{equation}
H=\tfrac{1}{\sqrt2}\begin{pmatrix}1&1\\1&-1\end{pmatrix},\quad
R_z(\theta)=\begin{pmatrix}e^{-i\theta/2}&0\\0&e^{i\theta/2}\end{pmatrix},\quad
R_x(\theta)=\begin{pmatrix}\cos\frac\theta2&-i\sin\frac\theta2\\-i\sin\frac\theta2&\cos\frac\theta2\end{pmatrix},
\end{equation}
and the two-qubit $\mathrm{CNOT}\ket{a,b}=\ket{a,a\oplus b}$. Measuring in the computational basis
returns $z$ with probability $|\braket{z}{\psi}|^2$ (the Born rule). A classical
\emph{statevector simulator} stores all $2^n$ complex amplitudes. With \code{complex64} this needs
$8\cdot2^n$ bytes: 2\,KB at 8 qubits, 512\,KB at 16 and 8\,MB at 20. A 58-qubit state would need
about 2.3 exabytes, which is why the features are first reduced with PCA. The paper's figure of
128\,GB for 58 qubits is a typo; $8\cdot 2^{58}$ bytes $\approx 2.3\times10^{18}$ bytes.

\paragraph{Useful identities.}
Writing $s(z)=1-2z\in\{+1,-1\}$ for the eigenvalue of $Z$ on bit $z$, we have
$R_z(\theta)\ket{z}=e^{-i\theta s(z)/2}\ket z$: $R_z$ is \emph{diagonal}. The sequence
$\mathrm{CNOT}_{ij}\,R_z^{(j)}(\theta)\,\mathrm{CNOT}_{ij}$ maps
$\ket{z_i,z_j}\mapsto e^{-i\theta s(z_i\oplus z_j)/2}\ket{z_i,z_j}$, and since $s(z_i\oplus z_j)=s(z_i)s(z_j)$,
\begin{equation}
\mathrm{CNOT}_{ij}\,R_z^{(j)}(\theta)\,\mathrm{CNOT}_{ij}=\exp\!\big(-i\tfrac{\theta}{2}Z_i Z_j\big),
\label{eq:zzgate}
\end{equation}
which is also diagonal. This is the ``ZZ interaction'' of the feature map, and it is the key to the
fast simulator in Section~\ref{sec:impcompute}.

\subsection{Quantum feature maps and kernels}
A \emph{feature map} is a parameterised circuit $U(\bm{x})$ that prepares $\ket{\phi(\bm x)}=U(\bm x)\ket{0^n}$.
The \emph{fidelity quantum kernel} is
\begin{equation}
K(\bm a,\bm b)=\big|\braket{\phi(\bm a)}{\phi(\bm b)}\big|^2=\big|\bra{0^n}U^\dagger(\bm a)U(\bm b)\ket{0^n}\big|^2 .
\label{eq:qkernel}
\end{equation}
The second form shows how hardware computes it: prepare $U^\dagger(\bm a)U(\bm b)\ket{0^n}$,
measure all qubits, and record how often the all-zeros string appears. $K$ is a valid PSD kernel,
because it is the inner product of the ``density-matrix features''
$\rho(\bm x)=\ket{\phi(\bm x)}\!\bra{\phi(\bm x)}$ under the Hilbert--Schmidt inner product~\cite{schuld2019,schuld2021}.
It also satisfies $K(\bm x,\bm x)=1$.

\paragraph{Separable $R_x$ encoding (Method 1).} With $\ket{\phi(\bm x)}=\bigotimes_i R_x(x_i)\ket0$,
the overlap factorises qubit by qubit. Since $R_x(a)^\dagger R_x(b)=R_x(b-a)$ and
$\bra0R_x(\theta)\ket0=\cos(\theta/2)$,
\begin{equation}
K_{R_x}(\bm a,\bm b)=\prod_{i=1}^n\cos^2\!\Big(\frac{b_i-a_i}{2}\Big).
\label{eq:rxkernel}
\end{equation}
This is a classical, cheap-to-evaluate kernel. It is a product of one-dimensional kernels, and
it closely resembles a periodic RBF kernel. No quantum computer is needed to evaluate it.

\paragraph{ZZ feature map (Method 3).} Following~\cite{havlicek2019}, one repetition is
\begin{equation}
U_{\text{ZZ}}(\bm x)=\underbrace{\prod_{(i,j)\in E}\exp\!\big(-i(\pi-x_i)(\pi-x_j)Z_iZ_j\big)\prod_i e^{-ix_iZ_i}}_{D(\bm x)\ \text{(diagonal)}}\;H^{\otimes n},
\label{eq:zzmap}
\end{equation}
where $E$ is the set of entangled pairs (adjacent pairs, \emph{linear}, in the original code) and
the whole block is repeated $r$ times. The code implements $e^{-ix_iZ_i}$ as $R_z(2x_i)$ and the
pair term as $\mathrm{CNOT}\cdot R_z(2(\pi-x_i)(\pi-x_j))\cdot\mathrm{CNOT}$, using
Eq.~\eqref{eq:zzgate}. Figure~\ref{fig:circuit} shows one repetition on four qubits. For $r\ge2$,
the Hadamard layers sandwiched between diagonal layers are what is conjectured to make the kernel
hard to compute classically~\cite{havlicek2019}.

\begin{figure}[H]
\centering
\begin{tikzpicture}[x=1cm,y=0.75cm,
  g/.style={draw,fill=white,minimum width=0.75cm,minimum height=0.5cm,font=\scriptsize},
  zz/.style={draw,fill=qblue!10,minimum width=1.25cm,font=\scriptsize}]
\foreach \q in {0,...,3}{
  \draw (-0.6,-\q) node[left]{\small$\ket{0}$} -- (11.3,-\q);
  \node[g] at (0.3,-\q) {$H$};
  \node[g] at (1.6,-\q) {$R_z(2x_\q)$};
}
\foreach \a/\xpos in {0/3.3,1/6.0,2/8.7}{
  \pgfmathtruncatemacro{\b}{\a+1}
  \fill (\xpos,-\a) circle (2.2pt);
  \draw (\xpos,-\a) -- (\xpos,-\b-0.15);
  \draw (\xpos,-\b) circle (5pt); \draw (\xpos-0.17,-\b)--(\xpos+0.17,-\b);
  \node[zz,minimum height=0.5cm] at (\xpos+1.0,-\b) {$R_z(2\bar x_\a\bar x_\b)$};
  \fill (\xpos+2.0,-\a) circle (2.2pt);
  \draw (\xpos+2.0,-\a) -- (\xpos+2.0,-\b-0.15);
  \draw (\xpos+2.0,-\b) circle (5pt); \draw (\xpos+1.83,-\b)--(\xpos+2.17,-\b);
}
\draw[decorate,thick,qgrey] (-0.1,0.55) -- (11.2,0.55);
\node[font=\scriptsize,qgrey] at (5.5,0.85) {one repetition (repeated $r$ times); $\bar x_i=\pi-x_i$};
\end{tikzpicture}
\caption{One repetition of the ZZ feature map with linear entanglement on 4 qubits. Every gate
except $H$ is diagonal in the computational basis.}
\label{fig:circuit}
\end{figure}

\subsection{Kernel concentration and bandwidth}
\label{sec:conc}
Quantum kernels have a characteristic failure mode called \emph{exponential concentration}~\cite{thanasilp2024}.
For expressive feature maps on many qubits, the states of two different inputs are nearly
orthogonal in the $2^n$-dimensional space, so $K(\bm a,\bm b)\approx 0$ for all $\bm a\ne\bm b$
and the Gram matrix approaches the identity. An SVM trained on $K\approx I$ can fit the training
set perfectly, but its decision function on new points is close to a constant. This is the
quantum analogue of an RBF kernel with $\gamma\to\infty$. On hardware it is worse, because
resolving values of order $2^{-n}$ needs exponentially many shots.

Shaydulin and Wild~\cite{shaydulin2022} and Canatar et al.~\cite{canatar2022} showed that
\emph{scaling the input}, $\bm x\mapsto\alpha\bm x$ with $\alpha<1$, acts as a bandwidth: it
moves the states closer together, restores off-diagonal mass, and can turn a kernel that does not
generalise into one that does. The original project had no such parameter. Every feature used the
full $[0,\pi]$ range ($\alpha=1$).

\subsection{When can a quantum kernel help?}
\label{sec:gd}
A quantum kernel can only beat classical methods if (a)~it cannot be computed efficiently
classically, \emph{and} (b)~it fits the data better than the classical kernels one would
actually use. Huang et al.~\cite{huang2021} made (b) testable with the \emph{geometric difference}
\begin{equation}
g(K_C\|K_Q)=\sqrt{\big\|\sqrt{K_Q}\,(K_C+\lambda I)^{-1}\sqrt{K_Q}\big\|_\infty},
\label{eq:gd}
\end{equation}
computed on the training Gram matrices, both normalised to trace $N$. If $g$ is small (close to 1),
then for \emph{every} labelling the classical kernel generalises about as well as the quantum one,
so no advantage is possible on this data. A potential advantage requires $g\gtrsim\sqrt N$. We
compute $g$ against the tuned RBF kernel for every run (Section~\ref{sec:evalsw}). Recent
large-scale benchmarking~\cite{bowles2024} also found that, under careful tuning, out-of-the-box
classical models usually match or beat quantum models on small tasks. That is one more reason to
give the classical baseline the same tuning budget as the quantum model.
